\begin{abstract}
Equivariant weight-space encoders are designed to exploit permutation symmetry in neural network weights, yet their sample efficiency advantage over plain encoders has never been measured in a controlled, shared-split study. We compare three conditions on the ModelZooDataset CIFAR-10 CNN zoo---a flat-MLP (plain), flat-MLP with permutation augmentation (PermAug), and GNN-NFN (structural equivariance)---across training set sizes \{100, 250, 500, 1000, full\}. GNN-NFN achieves a \textbf{sample efficiency advantage} over flat-MLP: it reaches 90\% of its peak R$^2$ at just $N{\approx}147$ training models, whereas flat-MLP requires the full dataset (${\sim}7{,}000$ models) to reach 90\% of its peak R$^2{=}0.886$---yielding an efficiency ratio of approximately \textbf{47$\times$}. Even at $N{=}1{,}000$, flat-MLP achieves only 84\% of its peak while GNN-NFN exceeds 97\%. However, this advantage is data-regime dependent: at $N{=}100$, PermAug (R$^2{=}0.138$) outperforms structural equivariance (GNN-NFN R$^2{=}{-}0.016$) (single-seed; multi-seed replication required), revealing a minimum-data threshold of approximately 150--250 models below which data augmentation of a plain MLP is superior. At full training scale, all encoders converge within 1\% R$^2$, consistent with expressivity equivalence theory. These findings provide a concrete decision boundary for practitioners choosing between encoder types under zoo-size constraints, and motivate a new research target: reducing the minimum-data threshold for equivariant graph encoders.
\end{abstract}
